\section{Mapping de la cadena industrial global de la inteligencia corporizada}

Esta sección organiza el mapa de la cadena industrial que presentó Xie Chen (谢晨). La inteligencia corporizada se divide en varios tipos de empresas porque debe inventar al mismo tiempo «la plataforma del vehículo de próxima generación» y «los algoritmos L4/L5 que corren sobre esa plataforma». La conducción autónoma cuenta al menos con una plataforma madura, el automóvil; en la inteligencia corporizada, en cambio, la plataforma, el cuerpo robótico, los sensores, las manos, los modelos, los datos y los escenarios están todos en una etapa temprana. Cuando la dificultad es tan alta, la cadena industrial tiende a dividir el trabajo.

En el programa se distinguen, en términos generales, cuatro tipos de empresas: empresas de hardware, empresas de Foundation Model, empresas que integran software y hardware para escenarios verticales, y empresas de infraestructura centradas en la simulación. Esta clasificación no es una etiqueta estática, sino que indica dónde coloca el riesgo cada empresa: las de hardware resuelven el cuerpo robótico y la cadena de suministro; las de modelos, el cerebro y la recipe de datos; las de implantación, el ROI del escenario; y las de simulación, los datos y el mundo de entrenamiento.

\begin{figure}[H]
\centering
\includegraphics[width=0.96\textwidth]{embodied-supply-chain-map.png}
\caption{Cadena industrial global de la inteligencia corporizada: división del trabajo entre empresas de hardware, de modelos base, de implantación vertical y de simulación. Diagrama conceptual propio, elaborado a partir de la conversación de 00:55:25--01:09:22.}
\end{figure}

\begin{knowledgebox}{Lectura de la figura: cuatro tipos de empresas corresponden a cuatro recursos escasos}
En la capa de hardware escasean un cuerpo robótico confiable, el costo y la cadena de suministro; en la capa de modelos escasean el cómputo, el talento y una recipe de datos de propósito general; en la capa de implantación vertical escasean los clientes, los escenarios y la operación; en la capa de simulación escasean un mundo físico de alta calidad, la cadena de herramientas y el ciclo cerrado de evaluación. La inteligencia corporizada es demasiado pesada, por lo que a corto plazo se parece más a un ecosistema industrial que a una sola empresa que lo abarque todo.
\end{knowledgebox}

\subsection{Hardware, modelos, implantación y simulación}

Esta subsección desarrolla uno por uno los cuatro tipos de empresas. Entre las empresas de hardware destacan Yushu (宇树) y las empresas de manos diestras, que llevaron el cuerpo robótico a la investigación y a la industria y, en cierta medida, establecieron un estándar de investigación. Las empresas de Foundation Model incluyen Physical Intelligence, Skild, NVIDIA/GEAR, Google/DeepMind, entre otras, que buscan el «cerebro» del robot y la scaling recipe. Las empresas de implantación que integran software y hardware incluyen Figure, Tesla Optimus, The Bot Company, DynaRobotics, entre otras, y se enfocan en escenarios verticales concretos. Las empresas de simulación, como Guanglun (光轮) y Genesis, construyen infraestructura en torno a Real2Sim, Simulation y Sim2Real.

La valoración que hace el ponente de Yushu resulta muy ilustrativa para la industria: ofrece el cuerpo robótico a la comunidad académica internacional, de modo que los artículos y los estudiantes adquieren el hábito de trabajar con robots de Yushu, y luego trasladan ese hábito a la industria. Esto muestra que el «estándar» en la capa de hardware no proviene solo de las especificaciones, sino también del ecosistema de desarrolladores, del ecosistema de publicaciones y de la movilidad del talento. En la capa de simulación, Guanglun aspira a ser ese mismo tipo de estándar: no un proveedor más, sino la base común de la simulación corporizada.

\begin{knowledgebox}{Tabla de roles de la cadena industrial}
\begin{tabularx}{\textwidth}{p{0.18\textwidth}p{0.32\textwidth}X}
\toprule
Rol & Tarea central & Principales dimensiones de competencia \\
\midrule
Empresas de hardware & Cuerpo robótico, manos, sensores y equipo de captura & Costo, confiabilidad, ecosistema de desarrolladores, cadena de suministro. \\
Empresas de modelos & Modelos base para robots, cerebro, VLA y scaling recipe & Cómputo, talento, pirámide de datos, capacidad de RL. \\
Empresas de implantación & Entrega para escenarios como fabricantes de automóviles, almacenes, hogares y restaurantes & ROI del cliente, integración de software y hardware, operación y seguridad. \\
Empresas de simulación & Generar mundos entrenables que sostengan el ciclo cerrado de datos y algoritmos & Activos físicos, solver, API, tasa de éxito de Sim2Real. \\
\bottomrule
\end{tabularx}
\end{knowledgebox}

\subsection{Diferencias de oportunidad entre China y Estados Unidos}

La subsección anterior trató la cadena industrial global; esta examina la estructura por país. Xie Chen considera que en Estados Unidos hay oportunidades de emprendimiento en la capa de modelos, porque su ecosistema pone más énfasis en la división industrial del trabajo y los usuarios finales tienen mayor disposición a pagar por software, inferencia y servicios de automatización. Empresas como OpenAI y NVIDIA pueden obtener ingresos muy altos mediante modelos, cómputo y ecosistema, por lo que la capa de modelos puede constituir empresas independientes.

La estructura de China es distinta. En el país, la disposición a pagar por software y por tokens es relativamente baja, mientras que las capacidades de hardware, cadena de suministro y equipo completo son más fuertes, de modo que muchas empresas siguen la ruta de integrar software y hardware, entregar soluciones de extremo a extremo y vender hardware. En el programa se menciona que empresas como Zijie (字节), Xiaomi (小米) y Lixiang (理想) podrían ser más adecuadas para desarrollar el «cerebro», porque cuentan con cómputo, talento en IA, puntos de acceso de producto, cadena de suministro o escenarios de vehículos o robots. Este juicio no significa que China carezca de capacidad en modelos, sino que el modelo de negocio y la organización industrial colocarán esa capacidad dentro de sistemas integrados más verticales.

\begin{figure}[H]
\centering
\includegraphics[width=0.94\textwidth]{china-us-embodied-opportunity.png}
\caption{Diferencias de oportunidad en inteligencia corporizada entre China y Estados Unidos: Estados Unidos tiene más oportunidades en la capa de modelos; China es más adecuada para integrar cerebro + cuerpo robótico + escenario. Diagrama conceptual propio, elaborado a partir de la conversación de 01:09:22--01:15:33.}
\end{figure}

\begin{knowledgebox}{Lectura de la figura: las diferencias provienen de la estructura de pago y de la división industrial del trabajo}
A la izquierda, en Estados Unidos es más fácil que surja una capa de modelos independiente, porque el pago por software, el pago por inferencia y la división del trabajo en el ecosistema permiten que las empresas de modelos obtengan ingresos. A la derecha, en China es más fácil integrar cerebro, cuerpo robótico y escenario, porque la cadena de suministro de hardware es fuerte, el cobro por software por separado es débil y las empresas de producto necesitan más un valor de extremo a extremo. Al leer la figura, no debe entenderse como una diferencia de nivel tecnológico, sino como una diferencia de estructura comercial.
\end{knowledgebox}

\subsection{La pirámide de datos y la convergencia de caminos distintos}

Esta subsección explica por qué rutas técnicas distintas podrían terminar convergiendo. En apariencia, algunos equipos dicen seguir la ruta de datos reales, otros la de simulación, unos ponen el énfasis en imitation learning y otros en RL. Sin embargo, Xie Chen observa que los clientes de primer nivel adoptan cada vez más, internamente, una data pyramid: en la base, datos de internet o de propósito general; en el medio, datos sintéticos y simulación; en la cima, una cantidad pequeña de datos reales de alta calidad.

Esta estructura piramidal concuerda con la lógica expuesta antes. Los datos de internet aportan conocimiento del mundo y preentrenamiento visual y de lenguaje; los datos sintéticos aportan entornos de tareas controlables, diversos y paralelizables; los datos reales aportan grounding, calibración y la validación final. Los equipos todavía difieren en RL, VLA, proporción de datos sintéticos y arquitectura de algoritmos, pero el «co-training real + sintético» ya es un consenso entre las instituciones principales.

\begin{importantbox}{Pirámide de datos}
Es probable que la inteligencia corporizada no dependa solo de datos reales ni solo de datos de simulación. La ruta más sólida es: en la base, datos de propósito general que construyen conocimiento del mundo; en el nivel medio, datos sintéticos o de simulación que amplían la distribución de tareas; en la cima, datos reales para calibración, demostración y validación del despliegue.
\end{importantbox}

\subsection{Resumen de la sección}

La cadena industrial de la inteligencia corporizada mantendrá la división del trabajo por mucho tiempo. El hardware, los modelos, la implantación y la simulación controlan cada uno recursos escasos distintos; las diferencias entre China y Estados Unidos provienen más del modelo de negocio y de la organización industrial; y las distintas rutas podrían converger finalmente en la pirámide de datos y en el entrenamiento mixto con datos reales y sintéticos.
